% Generated by 04_analyze.py. Do not edit by hand.
\newcommand{\FigOneCaption}{Evaluation protocol: held-out speed bands (PRIMARY protocol). Macro-F1 against the number of sensors $k$. Each subset is first averaged over the 10 splits; the box plots then summarise those per-subset means across all $\binom{7}{k}$ subsets of each size. Whiskers span the full minimum to maximum range rather than 1.5$\times$IQR: every subset is enumerated, so this is a complete population rather than a sample, and the extreme subsets are of direct interest. The oracle line is the maximum over those $\binom{7}{k}$ candidates selected using TEST scores and is an upper bound only. The 10 splits come from 2 repeats of 5-fold cross-validation over the same recordings and are therefore not mutually independent.}
\newcommand{\FigTwoCaption}{Evaluation protocol: held-out speed bands (PRIMARY protocol). Grouped permutation importance: the drop in inner-validation macro-F1 when all 18 features of one sensor are permuted together with a shared row permutation. Bars are the mean over the 10 splits and error bars their standard deviation. This measures each sensor's MARGINAL contribution, so a bar near zero may indicate redundancy given the other sensors rather than an absence of information. The 10 splits come from 2 repeats of 5-fold cross-validation over the same recordings and are therefore not mutually independent.}
\newcommand{\FigThreeCaption}{Evaluation protocol: held-out speed bands (PRIMARY protocol). Per-class recall against $k$ for the permutation-selected subsets. Predictions are pooled across all 10 splits before recall is computed. Several classes may sit at recall 1.0, where their lines coincide exactly; they are deliberately not offset or jittered. The 10 splits come from 2 repeats of 5-fold cross-validation over the same recordings and are therefore not mutually independent.}
\newcommand{\FigFourCaption}{Evaluation protocol: held-out speed bands (PRIMARY protocol). Missed-detection, false-alarm and misdiagnosis rates against $k$ for the permutation-selected subsets, computed from predictions pooled across all 10 splits. The 10 splits come from 2 repeats of 5-fold cross-validation over the same recordings and are therefore not mutually independent.}
\newcommand{\FigFiveCaption}{Evaluation protocol: held-out speed bands (PRIMARY protocol). Row-normalised confusion matrices for the full 7-sensor configuration and for the permutation-selected subset at $k=4$. Repeat 0 only, so every window is counted exactly once and the two panels are directly comparable.}
\newcommand{\FigFiveKOneCaption}{Evaluation protocol: held-out speed bands (PRIMARY protocol). Row-normalised confusion matrices for the full 7-sensor configuration and for the permutation-selected SINGLE sensor ($k=1$). This is the variant that shows where the configurations actually differ; the $k$ =4 variant compares two near-identical configurations by construction. Repeat 0 only, so every window is counted exactly once and the two panels are directly comparable.}
\newcommand{\FigSixCaption}{Evaluation protocol: held-out speed bands (PRIMARY protocol). Spearman correlation of the RMS feature between sensors, over all windows. The diagonal is 1.0 by definition and is masked. The quantity is signed, so the colour scale is symmetric about zero and derived from the off-diagonal data ($\pm0.90$), so the real structure is visible instead of being flattened by the diagonal.}
\newcommand{\FigSixBCaption}{Evaluation protocol: held-out speed bands (PRIMARY protocol). Mean absolute Spearman correlation between sensors, averaged over all 18 matched features rather than RMS alone, over all windows. The diagonal is masked. The quantity is non-negative, so the colour scale is sequential and runs from 0 to the largest off-diagonal value ($0.60$), rather than the symmetric scale used for the signed matrix. This is a broader redundancy measure than the RMS-only matrix; it is descriptive and enters no hypothesis test.}
